\documentclass[10pt,a4paper]{amsart}
\usepackage[margin=19mm]{geometry}
\usepackage{amsmath,amssymb,mathtools}
\usepackage[scr=rsfs,cal=euler]{mathalfa}
\usepackage{xfrac}
\usepackage[unicode,hidelinks,bookmarks=false]{hyperref}
\newcommand{\ie}{i.e.\ }
\newcommand{\cf}{cf.\ }
\newcommand{\resp}{respectively\ }
\newcommand{\Subsection}{\subsection}
\providecommand{\nobreakdash}{\nobreak\hbox{-}}
\setlength{\emergencystretch}{2em}
\multlinegap0pt
\usepackage{amssymb}
\usepackage{xcolor}
\usepackage{enumerate}
\usepackage{bm}
\usepackage[normalem]{ulem}
\usepackage{tikz}
\newtheorem{claim}{Claim}
\DeclareMathOperator{\fs}{FS}
\title{Folkman's theorem and the primes}
\author{David J. Fern\'andez-Bret\'on}
\date{Source: arXiv:2509.12025v3 (CC BY 4.0)}
\begin{document}
\maketitle

\noindent\emph{Source and licence.} Folkman's theorem and the primes. Version arXiv:2509.12025v3 (CC BY 4.0), distributed by arXiv under the Creative Commons Attribution 4.0 International licence, http://creativecommons.org/licenses/by/4.0/, as recorded in the arXiv licence metadata for this exact version and shown on the abstract page https://arxiv.org/abs/2509.12025v3. Full licence notice: \texttt{licenses/arxiv-2509.12025-CC-BY-4.0.txt}.\par\smallskip

\noindent\textit{Editorial scope.} Counted unit: the article's claim claim (source0.tex lines 214--216). The article's complete original proof of it is delivered with it (source0.tex lines 218--228, one \texttt{proof} environment of 11 lines). No mathematics is written, repaired or re-derived: the statement and the proof are the article's own lines, in the article's own notation, and no retained line is substituted or re-flowed.  No further source-local context is needed: every cross-reference of every retained line resolves inside the two delivered spans.  Nothing was omitted from any retained span, so the package carries no omission record.  No retained line contains a citation, so the wrapper carries no bibliography and no bibliography entry.  No retained line contains a figure environment, an image-inclusion command or drawing code, so no image is delivered and the package declares no asset dependency.  Editorial formatting only, in addition: an AMS \texttt{amsart} preamble that restates the article's own declarations and macros used by the retained lines, namely the environments \texttt{claim} on separate counters, together with the standard packages those lines need.  The retained environments print in this document as Claim 1 in order of appearance, whereas the article prints the same items with the numbering of its own full document.  The complete native source of the article as distributed in the arXiv e-print for this exact version is bundled unchanged as \texttt{sources/185077/source0.tex}.
\begin{claim}
For each $p\in\mathbb P$ and for each $\alpha$, there are no more than $p^4$ elements $a\in X$ such that $\nu_p(a)=\alpha$. 
\end{claim}

\begin{proof}[Proof of claim]
Suppose, seeking a contradiction, that there are more than $p^4$ such elements. Since $\xi_p(a)\mod p^2$ is one of $p^2$ possibilities, for all $a\in X$, by the pigeonhole principle one can find $p^2$ distinct elements $a_1,a_2,\ldots,a_{p^2}\in X$, along with some $0<A<p^2$, such that $\xi_p(a_i)\equiv A\mod p^2$ and $\nu_p(a_i)=\alpha$ for all $i\leq p^2$. Since $(A,p)=1$, there exists some $t$, $0<t<p^2$, such that $tA\equiv p\mod p^2$. Letting $b=a_1+\cdots+a_t$ and $B=\xi_p(a_1)+\cdots+\xi_p(a_t)$, we get
\begin{equation*}
b=a_1+\cdots+a_t=p^\alpha(\xi_p(a_1)+\cdots+\xi_p(a_t))=p^\alpha B,
\end{equation*}
where
\begin{equation*}
B=\xi_p(a_1)+\cdots\xi_p(a_t)\equiv A+\cdots+A=tA\equiv p\mod p^2,
\end{equation*}
meaning $\nu_p(B)=1$. Therefore $\nu_p(b)=\alpha+1$ while $\nu_p(a_1)=\alpha$, contradicting the monochromaticity of $\fs(X)$ (since $c$ includes the information about the parity of $\nu_p$).
\end{proof}

\end{document}
